\section{Possible statuses and the scalar constraints}
\label{sec:status-span}

Call two statuses available when both can occur in the same key cell.
We first show that every available pair has overlap bounded below by a
constant chosen before the partition. This lets us solve the scalar
equations using only the available pairs.

A solution supplies a short witness list for each of the four cross
edges. If there is no solution, a linear functional vanishes on every
available generator and detects the target. At a fixed key cell we may
vary the status on either side while holding the other fixed. The
functional therefore predicts a bit exactly. Preparation either rules
out a frequent prediction of this kind or makes its value on the target
zero.

Write the endpoints of the two units as $(A,1),(A,2)$ and
$(B,1),(B,2)$. In this section $h_{A,i}$ or $h_{B,j}$ denotes an atom's
role bit $a(\mathrm{atom})$. The basis status
vector at an endpoint is denoted by $\mathbf x$, and its dimension is
the recorded dimension of that endpoint's effective space $C_i$.

\subsection{How small can a positive overlap be?}

Fix a common partition given by at most \(R\) linear rows. Its cells
must be the full fibers of those rows. Discard the prepared rank exception
from \eqref{eq:prepared-rank}, for the joint row space at both endpoints.
This space has dimension at most \(2gR\). The exception has mixed mass
at most \(16\cdot2^{2gR+b-31n/32}\), before either leaf is chosen.

Every full unary status is a function of at most \(m\) affine linear
forms of the nominal input. Here \(m\) depends only on the construction
and the pin bound. Indeed each status bit has a matrix representative
of bounded rank. A rank \(j\) matrix factors through at most \(2j\)
linear forms in its quadratic expression. Collect these forms over the
finitely many bits. Include the tester bits, role, effective-basis values, and
internal bit \(p_i\).

For a good orientation, use the independent coordinates \((c,y,z')\)
from \Cref{sec:preparation}. Write the full status as
\(F(Ac+Bw+d)\), with \(w=(y,z')\). Conditional on \(c\), the affine
argument is uniform on \(Ac+d+\operatorname{im}B\). Its law depends
only on the image of \(Ac+d\) in the quotient by \(\operatorname{im}B\).
Consequently each conditional status probability \(h_v(c)\) depends
on at most \(m\) linear rows and satisfies
\[
 h_v(c)>0\quad\Longrightarrow\quad h_v(c)\ge2^{-m}.
\]
Include the diagonal bit in the full status before conditioning on it.
Since that bit is independent of \(c\), the same lower bound holds
for every positive probability after conditioning.

If two status supports meet at \(c\), take all the rows on which
either support depends. Every point with the same values as \(c\) on
those rows lies in both supports. This fiber uses at most \(2m\) rows. Its uniform mass is at least \(2^{-2m}\). The
reference cell law is within a factor of 2 of the uniform coordinate law,
by the Gram Fourier estimate. Thus
\begin{equation}\label{eq:status-support-gap}
 \sum_c\nu(c)h_v(c)h_w(c)=0
 \quad\text{or}\quad
 \sum_c\nu(c)h_v(c)h_w(c)\ge2^{-4m-1}.
\end{equation}
Choose \(0<\eta\le2^{-4m-1}\), also satisfying the positivity bounds
below. This choice precedes the common partition. These statements concern
inputs before the test for independence modulo \(F_0\). Equation
\eqref{eq:base-rejection} accounts for the subsequent rejection.

\subsection{Generators and the coupled linear system}

Use the finite tables from \Cref{sec:finite-partitions}, outside the unary
concentration and positivity exceptions. Labels may be excluded at each
endpoint. The excluded set includes the labels in its protected selector
space from \eqref{eq:label-budget}.

% lax begin Lax793645.StatusGenerators
\begin{definition}[Available generator]
\label{def:status-generator}
On the cross edge $(i,j)$, an available generator consists of a vector
\begin{equation}
\label{eq:status-generator-vector}
g=(q,h_{A,i},h_{B,j},\mathbf x_{A,i},\mathbf x_{B,j},p_{A,i},p_{B,j})
\end{equation}
together with data witnessing it. Specify a common tag and a surviving
label at each end. The full statuses, including the individual tester
values, must be attainable by generic inputs in a common cell. We call
these data a \emph{certificate} for \(g\). By
\eqref{eq:status-support-gap}, this is equivalent to their conditional
probability overlap being at least the fixed threshold $\eta>0$.
Their support overlap \eqref{eq:table-overlap} is at least as large.
Different certificates can give the same vector \(g\). We keep the
labels with each certificate when choosing lists, but take spans in
the vector space of the numerical entries.

Let $V_{ij}$ be the binary span of these generators. Its \emph{base
projection} forgets the final two $p$ entries and has target
\[
B_{ij}=\F\oplus\F\oplus\F\oplus
             \F^{\dim C_{A,i}}\oplus\F^{\dim C_{B,j}}.
\]
\end{definition}
% lax end

Let \(B_{\rm stat}\) bound the number of pairs of full numerical
statuses at any two endpoints. Fix feasible values of the diagonal bit
and a generic label on each side. On reference cells of sufficient mass
each prescribed role and basis assignment has conditional probability
at least \(\beta>0\). The cells omitted by the unary estimate have
total mass less than \(1/2\). Summing the full status overlaps compatible
with the two assignments therefore gives at least \(\beta^2/2\).
There are at most \(B_{\rm stat}\) summands. Take
\(B_{\rm stat}\eta<\beta^2/2\). Some certificate then gives the
prescribed assignment with overlap at least \(\eta\), for any surviving
pair of labels.

Consequently
\begin{equation}\label{eq:status-full-projection}
 V_{ij}\longrightarrow B_{ij}\text{ is surjective},\qquad
 \dim\ker(V_{ij}\longrightarrow B_{ij})\le2.
\end{equation}
This remains true after any exclusions leaving generic labels.

Fix a table of cross pairings. It prescribes the values on each
effective basis through the two functionals
\[
(a+u_{B,j}U_{A,i})|_{C_{A,i}},\qquad
(a+u_{A,i}U_{B,j})|_{C_{B,j}}.
\]
Write their coordinate vectors in the recorded bases as
$\mathbf t_{A,i;j}$ and $\mathbf t_{B,j;i}$, respectively. We seek one
sum vector $v_{ij}\in V_{ij}$ for each cross edge, subject to
\begin{equation}
\label{eq:status-local-targets}
\begin{aligned}
q_{ij}&=1,& h_{A,i;ij}+h_{B,j;ij}&=1,\\
\mathbf x_{A,i;ij}&=\mathbf t_{A,i;j},&
\mathbf x_{B,j;ij}&=\mathbf t_{B,j;i},
\end{aligned}
\end{equation}
\begin{equation}
\label{eq:status-coupled-targets}
\begin{aligned}
\sum_{i=1}^2(p_{A,i;ij}+h_{A,i;ij})&=0 &&(j=1,2),\\
\sum_{j=1}^2(p_{B,j;ij}+h_{B,j;ij})&=0 &&(i=1,2).
\end{aligned}
\end{equation}
These are the conditions \eqref{eq:recipe-scalars}--\eqref{eq:recipe-balance}.
For example, \(p_{A,i;ij}\) denotes the sum of the internal bits
at endpoint \(A_i\) over the list for edge \(ij\). This lets the two
lists at \(A_i\) have different sums.
Once the vectors \(v_{ij}\) are chosen, Lemma~\ref{lem:small-tables}
prescribes the remaining values on selected atoms. With matching keys,
Theorem~\ref{thm:gradient-realization} then completes the functionals.

% lax begin Lax793645.ScalarDuality
\begin{lemma}[Dual relations]
\label{lem:status-duality}
Any obstruction to the scalar system has a nonzero vector of weights
$(s_1,s_2,r_1,r_2)\in\F^4$. On every edge span $V_{ij}$ these give
a relation
\begin{equation}
\label{eq:status-dual-relation}
s_j(p_{A,i}+h_{A,i})+r_i(p_{B,j}+h_{B,j})
=\mathbf c_{ij}\cdot\mathbf x_{A,i}
 +\mathbf c'_{ij}\cdot\mathbf x_{B,j}
 +\beta_{ij}(h_{A,i}+h_{B,j})+\gamma_{ij}q
\end{equation}
for suitable coefficients. The scalar system is soluble if, for every
such relation, the sum of the right sides evaluated at the local targets
is zero.
\end{lemma}
% lax end

% lax begin Lax793645Proofs.Manuscript.ScalarDuality.counter_score
\begin{proof}
Apply finite-dimensional duality to the linear constraint map on
$\bigoplus_{i,j}V_{ij}$. The coefficients of the four coupled equations
are $s_j,r_i$, giving \eqref{eq:status-dual-relation} separately on each
edge. If they were all zero, the local projection onto
$(q,h_{A,i}+h_{B,j},\mathbf x_{A,i},\mathbf x_{B,j})$ would be
surjective by \eqref{eq:status-full-projection}, forcing all other dual
coefficients to vanish. Hence a nonzero obstruction uses some coupled
equation. Conversely the constraint target lies in the image exactly
when it annihilates every dual relation. The coupled targets are zero,
so their evaluated contribution is the sum stated in the lemma.
\end{proof}
% lax end

\subsection{Which predictions can a dual relation give?}

Say that a pair has a \emph{dual relation} if some permitted exclusions
and some nonzero coupled coefficients give \eqref{eq:status-dual-relation}
on every available generator. Each excluded set has size at most
$B_{\rm exc}$ and contains the mandatory exclusions. Such a relation
need not obstruct a particular target. Excluding all such relations will
therefore prove solvability for every table.

By \eqref{eq:status-support-gap}, the relation holds on every pair of
statuses that can occur in the same cell. For a fixed \(q,c\), both
conditional supports are nonempty. Every pair in their Cartesian product
is available. A relation of the form
\(a(v_A)+b(v_B)=\gamma q\) therefore forces each unary parity to be
constant on its support. This gives a common function of \(q,c\).
Each input point has positive mass in its cell, so the identity holds
at every input point.

% lax begin Lax793645.SevenPredictions
\begin{lemma}[Seven common predictor pairs]\label{lem:finite-dictionary}
Fix a second unit outside the rank and unary exceptions. There are at
most 7 pairs of functions of the common key cells such that, for each
first unit forming a good pair with it, the following holds. If there
is a dual relation with some \(s_j=1\), one pair of functions predicts
both of the first unit's endpoint bits exactly. The coefficients and
label exclusions may depend on the first unit.

\end{lemma}
% lax end
% lax begin Lax793645Proofs.Manuscript.SevenPredictions.exact_seven_alternative
\begin{proof}
The exact sections of \Cref{sec:preparation} put every endpoint prediction
in the form
\begin{equation}\label{eq:status-extracted-prediction}
 p_{A,i}=T(\lambda_{A,i},\mathrm{atom})+\alpha_{A,i}h_{A,i}
                 +Q_i(c).
\end{equation}
Use the later pin bound \(K\) here. The corresponding tester-room
inequality is recorded in Appendix~\ref{sec:parameters}.

At a fixed endpoint there is at most one such expression. To see this,
subtract two of them on a generic label surviving both exclusions.
On a cell where every role and basis assignment has positive
conditional probability, vary one bit while keeping the cell fixed.
The two expressions then have the same coefficient at that bit.
Doing this for each bit identifies \(\lambda\) and \(\alpha\). Once
the coefficients agree, evaluating at any input of a cell identifies
the value of \(Q\) there. Every reference cell has such an input by
\eqref{eq:prepared-rank}. Full freedom of the role and basis bits is
needed only for the first step.

For \(s_j=1\), rearranging \eqref{eq:status-dual-relation} gives
\begin{align}
 p_{A,i}+\mathbf c_{ij}\cdot\mathbf x_{A,i}
              +(1+\beta_{ij})h_{A,i}
 ={}&r_i p_{B,j}+\mathbf c'_{ij}\cdot\mathbf x_{B,j}
              +(r_i+\beta_{ij})h_{B,j}+\gamma_{ij}q.
 \label{eq:status-deterministic-comparison}
\end{align}
If \(r_i=0\), positivity on the second side forces
\(\mathbf c'_{ij}=0\) and \(\beta_{ij}=0\). The first side then has
the inactive predictor \(Q=0,\alpha=1\). If \(r_i=1\), both endpoints
have their unique predictors. They have the same \(\alpha\) and \(Q\).
Both diagonal reference laws give positive mass to the same common
cells. The two normalized functions are independent of \(q\), so
evaluating at the two values of \(q\) forces its coefficient to be zero.

Fix \(j\) with \(s_j=1\). Each first endpoint has predictor 0 or
\(Q_{B,j}\), according as \(r_i=0\) or \(1\). Allowing either choice
of \(j\) gives at most 7 possibilities for the first unit, namely
\[
 (0,0),\quad(0,Q_{B,1}),\quad(Q_{B,1},0),\quad(Q_{B,1},Q_{B,1}),
\]
\[
 (0,Q_{B,2}),\quad(Q_{B,2},0),\quad(Q_{B,2},Q_{B,2}).
\]
Discard any choice requiring a nonexistent endpoint predictor.
\end{proof}
% lax end

% lax begin Lax793645.RareObstructions
\begin{lemma}[No frequent obstruction without prediction]
\label{lem:status-global-prediction}
In the alternative without exact common predictions, the mass of good
independent unit pairs having a dual relation is at most \(.002\).
\end{lemma}
% lax end
% lax begin Lax793645Proofs.Manuscript.RareObstructions.witnessed_cost
\begin{proof}
Fix the second unit. Each pair of functions in
Lemma~\ref{lem:finite-dictionary} predicts on first units of mass less
than \(10^{-6}\), by the definition of the alternative without predictions. Their union has mass at most \(7\cdot10^{-6}\),
and hence at most \(.001\). Average over the independently drawn second
unit. Interchanging the units gives the same bound for relations with a
nonzero coefficient on the other side. Every nonzero coupled relation
belongs to one of these two events.

This bound concerns pairs outside the rank and unary exceptions.
Those exceptions have mass at most \(1/8\) for each unit, hence at most
\(1/4\) for the pair.

\end{proof}
% lax end

\subsection{Predicted scalar relations on four finite cells}

For each tag partition the keys by the pair \((f_1,f_2)\in\F^2\).
Use every attained pair. Each has reference mass at least
\(\eta=2^{-2R_*}/2\) on both diagonal slices, by the coordinate count
in Proposition~\ref{prop:status-preparation}. When both endpoints are inactive
the only pair is \((0,0)\).
Put \(e_i=0\) for an inactive endpoint,
and associate a distinct basis vector to each distinct active class.
Thus \(e_i=e_j\ne0\) precisely when both endpoints are active with
the same \(\alpha\) bit and the same predictor polynomials.

At any retained cell the role and basis coordinates can be prescribed
arbitrarily. Define
\emph{predicted generators} by
\[
 p_{A,i}=\boldsymbol\lambda_{A,i}\cdot\mathbf x_{A,i}
             +\alpha_i h_{A,i}+f_i,\qquad
 p_{B,j}=\boldsymbol\lambda_{B,j}\cdot\mathbf x_{B,j}
             +\alpha_j h_{B,j}+f_j.
\]
These are numerical generators. Their realization by successful queries,
using the exact prediction identities, is proved in \Cref{sec:rare-overlap}.
The table contractions are
\begin{equation}\label{eq:status-table-delta}
 \mathsf u_{B,j}(\Delta_A)=\sum_iu_{B,j}U_{A,i}(\lambda_{A,i}),\qquad
 \mathsf u_{A,i}(\Delta_B)=\sum_ju_{A,i}U_{B,j}(\lambda_{B,j}).
\end{equation}

% lax begin Lax793645.PredictedScalarSystem
\begin{lemma}[Predicted scalar compatibility]\label{lem:status-predicted}
For every compatible table satisfying \eqref{eq:endpoint-injection},
the scalar system has a solution in the spans of the predicted
generators. On each edge the solution is a sum of 1 to 3 such
generators with fresh generic labels.
\end{lemma}
% lax end
% lax begin Lax793645Proofs.Manuscript.PredictedScalarSystem.one_attained_scalar_soluble
\begin{proof}
Apply Lemma~\ref{lem:status-duality} to the spans of these generators.
Substitution and arbitrary choice of the free role and basis coordinates
force the exact coefficient identities
\begin{equation}\label{eq:status-forced-coefficients}
 \mathbf c_{ij}=s_j\boldsymbol\lambda_{A,i},\quad
 \mathbf c'_{ij}=r_i\boldsymbol\lambda_{B,j},\quad
 \beta_{ij}=s_j(1+\alpha_i)=r_i(1+\alpha_j).
\end{equation}
The remaining relation \(s_jf_i+r_if_j=\gamma_{ij}q\) holds on every
attained cell. Each pair of predictor values occurs on both diagonal
slices. Keeping these values fixed and taking \(q=0\) and \(q=1\) gives \(\gamma_{ij}=0\).
The remaining equality holds at every key. Together with
\eqref{eq:status-forced-coefficients}, this gives \(s_je_i=r_ie_j\).
If both endpoints are inactive, their exact identities and table compatibility annul
all coupled targets. Otherwise a nonzero weight pattern requires the
nonzero classes to coincide. In that case the weights are
\(s_i=r_i=\1_{i\in S}\) for their common support \(S\).
The marked role contributions cancel between units. The remaining
contribution is
\[
 |S|(1+\alpha)+\sum_{j\in S}\mathsf u_{B,j}(\Delta_A)
                  +\sum_{i\in S}\mathsf u_{A,i}(\Delta_B)=0
\]
by the compatibility table in \Cref{sec:preparation}. Thus every dual obstruction vanishes.

Subtract the predicted linear terms from the two final coordinates.
The constrained part of each generator is just \((q,f_i,f_j)\in\F^3\).
A target in their span is a sum of at most 3 available triples.
Its \(q\)-coordinate is one, so the sum is nonempty. Put all required
free role and basis coordinates on one occurrence, and zero on the
others. Every retained pattern allows those choices at arbitrary generic
labels. Choose labels greedily outside the pin exclusions, prediction
exceptions, and already chosen labels. There are at most 12 positions,
so \eqref{eq:selector-size} leaves a permissible choice at every step.
\end{proof}
% lax end

% lax begin Lax793645.SolublePairs
\begin{lemma}[Scalar solvability on positive mass]\label{lem:status-span}
In the alternative without common predictions of
Proposition~\ref{prop:status-preparation}, choose \(\eta>0\) below the
bounds in \eqref{eq:status-support-gap} and the proof of
\eqref{eq:status-full-projection}. On a set of pairs of mass bounded
below by a positive constant, the table is safe and consistent and
the scalar system is soluble after every permitted label exclusion.

\end{lemma}
% lax end
% lax begin Lax793645Proofs.Manuscript.SolublePairs.no_out
\begin{proof}
By Proposition~\ref{prop:status-preparation} the probability of a
safe pair is greater than \(.99\). Removing the rank and unary exceptions
costs at most \(1/4\), and Lemma~\ref{lem:status-global-prediction}
costs at most \(.002\) on the remaining pairs. A positive fraction
therefore survives all these restrictions. For any such pair,
Lemma~\ref{lem:status-duality} gives solvability for its consistent
table after every permitted choice of exclusions.

\end{proof}
% lax end

\subsection{Label deletion and fresh short lists}

It remains to replace a vector in an edge span by a short list of actual
generators, using distinct labels across both lists at an endpoint.
The base space may have large dimension, but every base vector has a
lift consisting of a single generator at fresh generic labels.
The discrepancy from the desired vector belongs to the kernel, which
has dimension at most 2 by \eqref{eq:status-full-projection}. We will
show that it has a basis whose vectors each use at most 3 generators.
This gives a list of length at most \(1+3+3=7\).

% lax begin Lax793645.LabelDeletion
\begin{lemma}[Deleting labels preserves the spans]\label{lem:span-continuation}
On a good pair of orientations, deleting fewer than one quarter of the
labels at either endpoint leaves every edge span unchanged, even if only
generic generators are retained.
\end{lemma}
% lax end
% lax begin Lax793645Proofs.Manuscript.LabelDeletion.generic_status_band_stable
\begin{proof}
Take a linear form that vanishes on the retained generic generators. Write it
as the sum of a unary parity from each endpoint, assigning its coefficient
of $q$ to one side. On each common cell, every pair of possible statuses
is available. Thus each unary parity is constant on its conditional
support, and the two sides give the same function of $q$ and the cell.

The exact section argument writes each such function as $Q(c)+\gamma q$,
with $Q$ quadratic. The resulting identity holds at every generic input
for each retained label. Setting the ordinary coordinates to zero gives
the shared-input identity as well. Fix the base input. Its two sides are polynomials
of degree at most 2 in the selector bits. Their difference vanishes
outside the deleted labels, so the Boolean minimum-weight bound makes it
zero at every label.

Every positive reference cell has positive empirical mass on a good
orientation. Hence the extended input identities hold on all conditional
status supports. The original linear form therefore vanishes on every
available generator, including those with deleted labels. The two spans
have the same annihilator and are equal.
\end{proof}
% lax end

% lax begin Lax793645.FreshWitnessLists
\begin{lemma}[Fresh lists of at most 7 generators]
\label{lem:fresh-lists}
On each pair with a soluble scalar system from Lemma~\ref{lem:status-span},
each scalar solution can be realized with 1 to 7 available generic
generators on each edge. Over both lists at an endpoint, the chosen
selectors are independent modulo its protected selector space. Each
generator retains its certificate, so the matched status overlap at
each position is at least \(\eta\).
\end{lemma}
% lax end

% lax begin Lax793645Proofs.Manuscript.FreshWitnessLists.seven_list
\begin{proof}
Let $d_S$ bound the dimension of every protected selector space.
Every selector space formed during the construction below has at most
$H_*=2^{d_S+34}$ labels. The parameter choice puts this below the
deletion allowance in Lemma~\ref{lem:span-continuation}. Fix a scalar solution
on the original edge spans. They remain unchanged as labels are reserved.

We prove a statement for one edge after any labels used on earlier
edges have also been reserved. Write its original span as $V$, its
base space as $B$, and its projection as $\pi$. Let $G_L$ be the
generators avoiding the current reserved sets $L$. Their span is still
$V$. Let \(H_L\) be spanned by the sums of at most 3 generators of \(G_L\)
that project to zero. Within each such sum the selector vectors must be
independent modulo the current protected spaces at both endpoints.
We claim

\begin{equation}
\label{eq:status-short-kernel}
H_L=\ker\pi.
\end{equation}

Take two generators lifting the same \(b\in B\). Choose a third
lift of \(b\) whose selectors are independent of those of either lift
modulo the protected spaces, avoiding \(L\) as well. Adding it to
either original lift gives an allowed sum of length 2 in \(H_L\).
The original lifts thus have the same class modulo \(H_L\). Denote
that class by \(\ell(b)\). A generator above zero already belongs to
\(H_L\), so \(\ell(0)=0\).

Choose lifts of \(b,c,b+c\) in turn, preserving independence of the
selectors at both endpoints. Their sum projects to zero and has length
3, so \(\ell(b)+\ell(c)=\ell(b+c)\). This makes \(\ell\) linear.
It is inverse to the projection \(V/H_L\to B\) on each generator.
The generators span \(V\), proving \eqref{eq:status-short-kernel}.

Now lift the base projection of the desired vector with one generator
and reserve its selectors. The discrepancy belongs to \(\ker\pi\).
As long as the kernel vectors chosen so far span a proper subspace
of \(\ker\pi\), choose a sum of at most 3 fresh generators lying
outside that subspace. Such a sum exists by
\eqref{eq:status-short-kernel}, since Lemma~\ref{lem:span-continuation}
preserves the edge span after each reservation. Reserve its selectors
before choosing another sum. In at most 2 steps the resulting sums
form a basis of \(\ker\pi\). Add those needed for the discrepancy
to the initial lift. Its list is nonempty and has length at most 7.
Because each reservation included all previous choices, the selected
labels are independent modulo the protected spaces at both endpoints.

Carry out this procedure on the four edges in turn, enlarging the
protected spaces by all earlier choices. A unit uses at most 28 selected
vectors. The auxiliary short-word comparisons require at most 6 more
vectors. Every excluded span therefore contains at most $2^{d_S+34}=H_*$
labels. The deletion lemma applies at every step. The generators retain their
certificates of overlap at least $\eta$, proving the lemma.
\end{proof}
% lax end

\subsection{A fixed lower bound for the finite score}

% lax begin Lax793645.UnpredictedOverlap
\begin{proposition}[Positive-mass overlap]\label{prop:positive-overlap}
In the alternative without common predictions of Proposition~\ref{prop:status-preparation},
there is a fixed option for which the expected overlap of the
accepted key measures is bounded below by a positive constant along
a subsequence. Restoring any permitted restriction of mass
\(2^{-o(N)}\) still gives the overlap required by
Proposition~\ref{prop:collision}.
\end{proposition}
% lax end
% lax begin Lax793645Proofs.Manuscript.UnpredictedOverlap.one_quiet_end
\begin{proof}
If the conclusion failed, all the finitely many $I_n(\omega)$ would tend
to zero. Choose $\eta$ as in Lemma~\ref{lem:status-span}. Let $c>0$ be a
lower bound for the mass of safe pairs with a soluble scalar system after the exceptional
losses. Apply Lemma~\ref{lem:fresh-lists}. Each list has at most 7 entries
and every full-status certificate has overlap at least $\eta$.
There are four edges, hence at most 28 factors. Choose $\eta\le1$ and put
\begin{equation}\label{eq:finite-score-lower}
 a=\kappa_{\min}^{\,2}\eta^{28}>0.
\end{equation}
The scalar and gradient consistency equations provide the required
binary prescriptions. With the safe table they give a usable option
of score at least $a$ for each of these pairs. Tester summaries were
included in the certificates at the outset, so no subsequent subdivision
reduces this bound.

Thus \(\PP\{\exists\omega:s_\omega\ge a\}\ge c\). Apply
Proposition~\ref{prop:finite-score-bound} with this \(a\) and
\(\epsilon=c/2\). It gives the opposite bound for large \(n\),
a contradiction. If a restriction has mass \(2^{-o(N)}\), restoring
it for two independent draws multiplies the lower bound by
\(2^{-o(N)}\), as in Proposition~\ref{prop:collision}.

\end{proof}
% lax end
